(1) 小扰动声压满足 \(p_{tt}=c^2p_{xx}\)。设 \(p(x,t)=P(x)e^{-i\omega t}\)，有 \(P''+k^2P=0\)，其中 \(k=\omega/c\)。封闭端法向速度为零，线性动量方程 \(\rho v_t=-p_x\) 给 \(P'(0)=0\)；开端在本近似下声压固定，\(P(L)=0\)。

(2) 由 \(P'(0)=0\)，一般空间解取 \(P=A\cos(kx)\)。开端条件 \(\cos(kL)=0\)，所以
\[
k_nL=(n+\tfrac12)\pi,\qquad
f_n=\frac{\omega_n}{2\pi}=\frac{(2n+1)c}{4L},\quad n=0,1,2,\ldots
\]
所有频率是基频 c/(4L) 的奇整数倍。

(3) 谐波动量方程 \(-i\omega\rho V=-P'\)，故速度复幅度 \(V=P'/(i\omega\rho)=ikA\sin(kx)/(\omega\rho)\)（整体相位随时间惯例而定）。因此封闭端 x=0 是速度节点、压力腹点，开放端 x=L 是压力节点、速度腹点。

(4) 当管长 \(L\mapsto2L\)，所有频率减半，不改变奇次谱结构。由于混合 Neumann–Dirichlet 边界只能容纳 \(1/4,3/4,\ldots\) 波长对应的空间驻波，不出现两端同类边界下的偶次谐波。
